\section{Quadratic transfer of the Case III tensor}
\label{sec:tensor-transfer66}

Fix an algebraic closure $\Omega$ of $\mathbb Q$, any prime $\ell$,
an isomorphism $\iota:\overline{\mathbb Q}_\ell\simeq\mathbb C$,
and an imaginary quadratic field $K'\subset\Omega$.  Write
$G_E=\operatorname{Gal}(\Omega/E)$ and
$\chi=\chi_\ell$, with geometric Frobenius of norm $q$ acting by
$q^{-1}$ on $\chi$.  Let the continuous rank-two representation $\rho$
of $G_{\mathbb Q}$ satisfy the source hypotheses
$\texttt{IsAutR}$, $\texttt{IrrOpen}$ and $\texttt{PR}$.
The last hypothesis supplies a regular algebraic cusp form $\pi$ over
$\mathbb Q$ attached to $\rho$. At every finite place its component
is spherical or has two unlinked characters as cuspidal support.
Use this same $\pi$ throughout.  The arithmetic inputs below are cyclic
base change and quadratic induction with local compatibility (SX-BC),
rank-two polarized CM local--global compatibility away from $\ell$
(SX-G), and the class field theory and local principal series
statements in SX-C.  We use strong multiplicity one for cuspidal
$\mathrm{GL}_n$.  The geometric attachment in rank $n$ multiplies
the usual Satake parameters by $q^{(n-1)/2}$.

% BEGIN F66A1
\paragraph{F66A1. A single exponent for both embeddings.}
Fix $g_0\notin G_{K'}$, $m\in\mathbb N$, a coefficient ring
$O_1$ and $\kappa:G_{K'}\to\overline{\mathbb Q}_\ell^\times$,
and suppose that a witness $\texttt{TauDat}$ exists for $(m,\kappa)$.
Its positivity clause gives $m\ge1$.
Quadratic base change of $\pi$ to $K'$ is cuspidal. Indeed failure
of cuspidality would give a self-twist by the nontrivial character
of $K'/\mathbb Q$. Comparison of good Frobenius polynomials,
Chebotarev and semisimplicity give the corresponding self-twist of
$\rho$. Its intertwiner commutes with $\rho(G_{K'})$, so is scalar
by open-subgroup irreducibility. On the other coset it would then
equal its negative. The base change $P_0$ is regular algebraic
and attaches $\rho|_{G_{K'}}$ by the archimedean and unramified
local formulas.

For either embedding $e:K'\hookrightarrow\mathbb C$, let $A_e$
denote the two entries of the infinitesimal character of $P_0$.
They are distinct and belong to $\tfrac12+\mathbb Z$. Form
\[
 S=\bigcup_{e:K'\hookrightarrow\mathbb C}
        \{u-v\mid u\in A_e,\ v\in A_{\bar e}\}.
\]
There are at most eight elements.  The nine numbers $2mj$ with
$1\le j\le9$ are distinct. Fix one $j$ with $2mj\notin S$.
For \emph{every} later $\texttt{TauDat}$ witness for
$(mj,\lambda)$, $\lambda=\kappa^j$, its algebraicity clause gives
a Hecke character $Y$ realizing $\lambda$ and an embedding $e$
with exponents $-mj,+mj$ at $e,\bar e$. Thus the two
infinitesimal characters of $P_0\otimes Y$ are
$A_e-mj$ and $A_{\bar e}+mj$. A common entry would give
$u-v=2mj$. The four entries are distinct and lie in
$\tfrac32+\mathbb Z$ as well. This choice of $j$ precedes
$\rho_1,\tau,s$, the Hecke-character witness and its orientation,
and the field $L$ below.
% END F66A1

% BEGIN F66A2
\paragraph{F66A2. The tensor over the quadratic compositum.}
Take any integral $\rho_1,\tau:G_{\mathbb Q}\to M_2(O_1)$ and
$s$ satisfying the source comparison $\rho_1\otimes
\overline{\mathbb Q}_\ell\simeq\rho$ and the chosen
$\texttt{TauDat}$.  Let $R=\mathrm{tenB}(\rho_1,\tau)$.
Take any totally complex abelian number field $L/\mathbb Q$
in $\Omega$ with $K'\not\subset L$, with $\texttt{GoodAt}(K',L,\rho)$
and $R|_{G_L}$ irreducible.  Put $M=LK'$.  Since $K'$ is quadratic,
$M/L$ is quadratic and $M/K'$ is abelian. The fields $M$ and $L$
are CM. Indeed complex conjugation is a nonidentity central
involution in the abelian Galois group and fixes a totally real
subfield of index two.

In the second tensor factor put $f_+=(1,1)$ and $f_-=(1,-1)$
over $\overline{\mathbb Q}_\ell$. On $G_{K'}$ the matrix of $\tau$
has rows $(a,b),(b,a)$ with $a^2-b^2=1$. Its eigencharacters
on these lines are $a+b=\lambda$ and $a-b=\lambda^{-1}$.
On the other coset the rows are $(a,b),(-b,-a)$, so the two
lines are exchanged. This calculation uses characteristic zero
also when $\ell=2$. In the signed tensor basis of Q65P1, the
original four matrices represent the ordinary tensor product.
The two planes obtained by tensoring with $f_+,f_-$ are exchanged
by every $x\in G_L\setminus G_M$. On the first plane $G_M$ acts by
$r=\rho|_{G_M}\otimes\lambda|_{G_M}$. Sending the identity coset
in the induced space to this plane and the $x$-coset to its image
gives an isomorphism
\[
             R|_{G_L}\simeq\operatorname{Ind}_{G_M}^{G_L}r.
\]
In particular at $h\in G_M$ its characteristic polynomial is the
product for $r(h)$ and $r(x^{-1}hx)$, and at $x$ it is
$\det(X^2 I_2-r(x^2))$. These follow from the diagonal and
off-diagonal block matrices, respectively. Twisting by $\lambda$
preserves irreducibility on every open subgroup of $G_{K'}$.
Continuity follows from the continuity of $\rho$ and of $\lambda$.

Let $P$ be the base change to $M$ of $P_0\otimes Y$. Twisting
attaches $\rho|_{G_{K'}}\otimes\lambda$. Factor $M/K'$ into cyclic
prime-degree layers. The self-twist argument of F66A1 applies
at every layer, using irreducibility of the restriction to the
next field. At unramified places of residue degree $f$, raising
the Satake parameters to the $f$th power commutes with the
rank-two factor $q^{1/2}$ because $q_w=q_v^f$. Thus $P$ is
cuspidal and attaches $r$, and its infinitesimal characters
are pulled back from $P_0\otimes Y$. Each complex embedding of
$L$ has two extensions to $M$. They induce distinct embeddings
of $K'$, since $L$ and $K'$ generate $M$. Their distinct
infinitesimal-character multisets from F66A1 in particular
give $P^x\not\simeq P$ for the nonidentity automorphism of $M/L$.
% END F66A2

% BEGIN F66A3
\paragraph{F66A3. One local window upstairs and sphericity at nonsplit places.}
First take a real norm twist of $\pi$ which is unitary. Such
a global twist exists because the absolute value of the central
character is a power of the idele norm. In this normalization
each finite component is a generic unitary principal series.
In the plain alternative this follows from its unlinked support,
and in the spherical alternative a nongeneric one-dimensional
constituent cannot be a local factor of a cusp form on
$\mathrm{GL}_2$.  Write $u_1,u_2$ for the two inducing characters.
Unitarity identifies their multiset with
$\{\bar u_1^{-1},\bar u_2^{-1}\}$. If they are separately fixed,
both have absolute value one. Otherwise they have equal restrictions
to units, since smooth characters of a compact group have finite
image. Twist off that common unitary character, extended by value
one on a uniformizer. The resulting spherical unitary component has
parameters $zv,z^{-1}v$ with $z\ge1$ and $|v|=1$. For its spherical
unit vector the normalized adjacency operator is $q^{-1/2}$ times
a sum of $q+1$ unitary translates. Therefore
$z+z^{-1}\le q^{1/2}+q^{-1/2}$. Equality would force all the
translates into the original line, since their sum is a nonzero
multiple of the vector. The maximal compact subgroup and its
uniformizer double coset generate the group. This would make the
irreducible representation one-dimensional, contradicting genericity.
In either case each absolute parameter is strictly between
$q^{-1/2}$ and $q^{1/2}$.

Undoing the global norm twist translates both exponents by
one common real number. Under local base change a character
is composed with the norm. The valuation of a uniformizer norm
is the residue degree $f$, and smooth characters have absolute
value one on units. Thus the exponent measured in base $q_w=q_v^f$
is unchanged. The Hecke character $Y$ translates the exponents
by a second common real number, since the absolute value of an
idele class character is trivial on the compact norm-one class
group and factors through the idele norm. Consequently there
is $t\in\mathbb R$ such that every finite component of $P$ is
toral, with
\[
   q_w^{t-1/2}<|\nu(\varpi_w)|<q_w^{t+1/2}
\]
for all its support characters $\nu$. No two such characters
have quotient $|\cdot|^{\pm1}$. In particular the local parameters
have zero monodromy. These computations use the local principal
series formulas of SX-BC at ramified places as well.

For $v$ of $L$ above a rational prime $p$, the first alternative
of $\texttt{GoodAt}$ says that every decomposition group over $p$
fixes $K'$. Its intersection with $G_L$ then fixes $M$, so $v$
splits in $M/L$. At any nonsplit $v$, the second alternative
therefore gives $p\ne\ell$, trivial image of the $L$-inertia in
$\operatorname{Gal}(M/L)$, and trivial action of $\rho$ on that
inertia. Thus $M_w/L_v$ is unramified of degree two. For $w$
above $v$, every element of its inertia lies over $K'$ and lies
in rational inertia away from $\ell$. By $\texttt{TauDat.unr}$
its $\tau$-matrix is the identity, so $\lambda$ is one there.
Both factors of $r$ are unramified at $w$.

Here are the polarization and compatibility required to infer
sphericity of $P_w$. Choose rational complex conjugation $c$.
It acts nontrivially on $K'$ and interchanges $f_+$ and $f_-$,
so $\lambda^c=\lambda^{-1}$. With
$J=\left(\begin{smallmatrix}0&1\\-1&0\end{smallmatrix}\right)$
and $B=J\rho(c)^{-1}$ one has
\[
    r(h)^t B r(chc^{-1})=\det\rho(h) B
                      \qquad(h\in G_M).
\]
This follows from $A^tJA=\det(A)J$. The multiplier is algebraic.
The algebraic central character $\omega_\pi$ has $\ell$-adic
character $\omega_{\pi,\ell}$, and
$\det\rho=\omega_{\pi,\ell}\chi^{-1}$ by comparing geometric
Frobenius at good primes and applying Chebotarev.
On the automorphic side, base change of $\pi$ is conjugation
invariant and its rank-two dual is its twist by the inverse central
character.  Since $Y^c=Y^{-1}$ by reciprocity and density, one gets
$P^c\simeq P^\vee\otimes(\omega_\pi\circ N_{M/\mathbb Q})$,
where characters act through determinant. This multiplier is pulled
back from $M^+$ and has uniform sign at its real places. Multiplying
the character on $M^+$ by the quadratic character of $M/M^+$
changes that sign at every real place and leaves its norm pullback
fixed. This gives the sign required in SX-G. Its semisimple Galois
representation is $r$ by good Frobenius comparison and irreducibility.
At $w\nmid\ell$ semisimplified local--global compatibility now
identifies the two local support characters on inertia with the
semisimple inertia characters of $r$ up to an unramified norm twist.
At our nonsplit place $r$ is trivial on inertia. Equivalently, for
every inertia element the two linear factors multiply to $(X-1)^2$,
so both characters are one there. The unlinked principal series
is therefore spherical. This proves precisely that at each nonsplit
$v$, $M/L$ is unramified and $P_w$ is spherical for every $w\mid v$.
% END F66A3

% BEGIN F66A4
\paragraph{F66A4. The two universal conclusions for the same tensor.}
By F66A2 and the quadratic cuspidality criterion of SX-BC,
$P$ induces to a cuspidal $P_4=\operatorname{AI}_{M/L}P$.
Its infinitesimal character at any embedding of $L$ is the sum
of the two multisets in F66A1, hence regular algebraic in rank
four. Put $\Pi=P_4\otimes|\det|$. The twist translates the
infinitesimal characters by an integer, preserving this property.
To check attachment, at a good split place the usual parameters
of $P_4$ concatenate the two upstairs pairs. At a good inert place
they are the two square roots, with both signs, of each upstairs
parameter. Upstairs multiplies the usual parameter by $q_w^{1/2}$.
Thus the geometric eigenvalues of $R$ equal $q_v^{1/2}$ times
the usual parameters of $P_4$ in both cases, by the two
characteristic-polynomial formulas of F66A2. Rank-four attachment
uses $q_v^{3/2}$, and $|\det|$ multiplies each usual parameter
by $q_v^{-1}$. This proves attachment to $\Pi$. The irreducibility
assumed for $R|_{G_L}$ supplies the semisimplicity clause.

For the all-place claim take $t$ from F66A3. At a split place
$P_4$ has the four upstairs support characters by the local
induction formula. At a nonsplit place the extension and both
upstairs characters are unramified by F66A3. A character with
value $b$ on the upstairs uniformizer induces to the two
unramified characters with values $\sqrt b,-\sqrt b$ downstairs.
Since $q_w=q_v^2$, both lie in the strict window with center $t$
in base $q_v$. The common window at split places holds with the
same center, also above $\ell$. For every pair of the four characters
their absolute ratio is strictly between $q_v^{-1}$ and $q_v$.
The unlinked-character criterion of SX-C identifies the local
representation with the irreducible principal series of these
characters, so its entire support is toral. Twisting by $|\det|$
changes the center to $t-1$ at every finite place.

Finally, take \emph{any} regular algebraic cuspidal representation
of $\mathrm{GL}_4(\mathbb A_L)$ attached to $R|_{G_L}$. Its Satake
multisets agree almost everywhere with those of $\Pi$, by attachment
to the same geometric characteristic polynomials. Strong multiplicity
one identifies it with $\Pi$. All its finite components thus have
the same strict window as above. This supplies the full universal
$\texttt{AllW}$ for $R|_{G_L}$. Since $L$ was arbitrary under the
source hypotheses, $\texttt{AutOffS}$ and $\texttt{OffW}$ both
hold for each tensor with the one exponent of F66A1. This proves
$\texttt{RankFourAutomorphyOfTensor}$ with the declared classical
inputs.
% END F66A4
